% Лекц 1: семинар ба гэрийн даалгавар. Бодлогууд хөнгөнөөс хүнд рүү эрэмбэлэгдсэн.

\ProblemSet{Семинар 1}
\begin{frame}[allowframebreaks]{Семинар 1}
\pgroup{Хэллэг ба логик холбоос.}
\begin{plist}
\pr{Дараах өгүүлбэрүүдийн аль нь хэллэг вэ? Хэллэг бол үнэнийг тогтоо: "7 нь тэгш"; $5<9$; "Хаалгаа хаа!"; $2+3=6$.}
\pr{$P$: "$3<5$", $Q$: "7 нь тэгш" бол $P\land Q$, $P\lor Q$, $\lnot P$, $P\Rightarrow Q$, $Q\Rightarrow P$ хэллэгүүдийн үнэн худлыг тогтоо.}
\pr{$x>2\Rightarrow x^2>4$ хэллэгийг монгол хэлээр бич.}
\pr{"Хэрэв тоо 6-д хуваагдвал 3-д хуваагдана" хэллэгийн урвуу хэллэгийг бичээд үнэн эсэхийг дүгнэ.}
\end{plist}
\pgroup{Логикийн хуулиуд, тавтологи ба зөрчил.}
\begin{plist}
\pr{Үнэний хүснэгтээр шалга: $(P\Rightarrow Q)\Leftrightarrow(\lnot P\lor Q)$.}
\pr{$F=P\land\lnot P$ хэллэг зөрчил болохыг батал.}
\pr{Де Морганы хуулийг хүснэгтээр батал: $\lnot(P\land Q)\Leftrightarrow(\lnot P\lor\lnot Q)$.}
\pr{Батал: $P\land(P\lor Q)\Leftrightarrow P$ (шингээх хууль).}
\pr{$(P\Rightarrow Q)\lor(Q\Rightarrow P)$ тавтологи мөн эсэхийг тогтоо.}
\pr{Үнэний хүснэгтээр шалга: $\big(P\land(Q\lor R)\big)\Leftrightarrow\big((P\land Q)\lor(P\land R)\big)$ (дистрибутив хууль).}
\pr{Эквивалент болохыг батал: $(P\Rightarrow Q)\Leftrightarrow(\lnot Q\Rightarrow\lnot P)$.}
\end{plist}
\pgroup{Баталгааны аргууд, квантор.}
\begin{plist}
\pr{Шууд батал: хэрэв $n$ сондгой бол $n^2$ сондгой.}
\pr{Шууд батал: хэрэв $a$ бүхэл бол $a^2+a$ тэгш.}
\pr{Эсрэгээс нь батал: хэрэв $ab$ сондгой бол $a$ ба $b$ сондгой.}
\pr{$P$: "$n$ тэгш", $Q$: "$n^2$ нь 4-т хуваагдана". $P\Rightarrow Q$ хэллэгийг $n=6$, $n=7$, $n=10$ үед шалга. $P\Leftrightarrow Q$ теоремыг батал.}
\pr{Дараах өгүүлбэрийн үгүйсгэлийг кванторыг зөв байршуулж бич: $\forall y\in\R,\ \exists x\in\R,\ (f(x)=y)$.}
\pr{$a_1,a_2,\ldots$ дарааллын хувьд $\exists M>0,\ \forall n\in\N,\ (|a_n|<M)$ чанар биелдэг байг. Ийм дарааллын жишээ гарга. Энэ чанарын үгүйсгэлийг бич.}
\pr{$\sqrt6$ иррационал болохыг батал.}
\pr{Индукцээр батал: $1^2+2^2+\cdots+n^2=\dfrac{n(n+1)(2n+1)}{6}$.}
\pr{$a,b\in\R$, $n\in\N$ бол $a^n-b^n=(a-b)\sum_{k=0}^{n-1}a^{n-1-k}b^k$ болохыг батал. Мөн $n$ сондгой үед $a^n+b^n=(a+b)\sum_{k=0}^{n-1}(-1)^ka^{n-1-k}b^k$ болохыг батал.}
\end{plist}
\end{frame}

\ProblemSet{Гэрийн даалгавар 1}
\begin{frame}[allowframebreaks]{Гэрийн даалгавар 1}
\pgroup{Хэллэг, логик холбоос, үгүйсгэл.}
\begin{plist}
\pr{Дараах өгүүлбэрүүдийн аль нь хэллэг вэ? "$x+1=3$"; "Улаанбаатар бол Монголын нийслэл"; "2 бол анхны тоо"; "Энэ өгүүлбэр худал".}
\pr{$P$: "$2^{10}>1000$", $Q$: "17 нь 3-т хуваагдана" бол $P\land Q$, $P\lor Q$, $P\Rightarrow Q$, $Q\Rightarrow P$, $\lnot P\lor Q$-ийн үнэнийг тогтоо.}
\pr{Үгүйсгэлийг бич: "Ангийн бүх оюутан даалгавраа хийсэн"; "1-ээс бага натурал тоо олдоно".}
\pr{"Хэрэв $x=2$ бол $x^2=4$" хэллэгийн урвуу, эсрэг, урвуу-эсрэг хэллэгийг бичиж, алинь нь үнэн болохыг тогтоо.}
\pr{Үнэний хүснэгтээр шалга: $(P\land Q)\Rightarrow P$.}
\pr{Үнэний хүснэгтээр шалга: $(P\lor Q)\Leftrightarrow(\lnot P\Rightarrow Q)$.}
\pr{Эквивалент болохыг батал: $\lnot(P\Rightarrow Q)\Leftrightarrow(P\land\lnot Q)$.}
\pr{$(P\Rightarrow Q)\Rightarrow(Q\Rightarrow P)$ тавтологи мөн үү?}
\end{plist}
\pgroup{Эсрэг жишээ, шууд ба эсрэгээс батлах.}
\begin{plist}
\pr{Эсрэг жишээ гарга: "Бүх тэгш тоо анхны тоо".}
\pr{Эсрэг жишээ гарга: "Бүх анхны тоо сондгой".}
\pr{Эсрэг жишээ гарга: "Бүх бүхэл тоо анхны эсвэл тэгш тоо".}
\pr{Эсрэг жишээ гарга: "Хэрэв $ab$ тэгш бол $a$ ба $b$ хоёулаа тэгш".}
\pr{Шууд батал: хэрэв $a,b$ тэгш бол $a+b$ тэгш.}
\pr{Эсрэгээс нь батал: хэрэв $n^2$ сондгой бол $n$ сондгой.}
\pr{Эсрэгээс нь батал: хэрэв $3\mid n^2$ бол $3\mid n$.}
\pr{Батал: хэрэв $a$ нь 4-т хуваагддаг бол $a^2$ нь 8-д хуваагдана.}
\pr{Тавтологи болохыг батал: $\big((P\Rightarrow Q)\land(Q\Rightarrow R)\big)\Rightarrow(P\Rightarrow R)$.}
\pr{Тавтологи болохыг батал: $\big((P\lor Q)\land(\lnot P\lor R)\land(\lnot Q\lor R)\big)\Rightarrow R$.}
\end{plist}
\pgroup{Квантор.}
\begin{plist}
\pr{$a_1,a_2,\ldots$ дарааллын хувьд $\forall M>0,\ \exists n\in\N,\ (a_n>M)$ чанар биелдэг байг. Ийм дарааллын жишээ гарга. Энэ чанарын үгүйсгэлийг бич.}
\pr{Үгүйсгэлийг кванторыг зөв байршуулж бич: $\forall\varepsilon>0,\ \exists\delta>0,\ \big(|x-1|<\delta\Rightarrow|f(x)-f(1)|<\varepsilon\big)$.}
\end{plist}
\pgroup{Математик индукц ба эсрэгээс батлах.}
\begin{plist}
\pr{Индукцээр батал: $1+2+4+\cdots+2^{n-1}=2^n-1$.}
\pr{Индукцээр батал: $7^n-1$ нь 6-д хуваагдана.}
\pr{Индукцээр батал: $n\ge5$ үед $2^n>n^2$.}
\pr{Бүх анхны $p$-ийн хувьд $p\mid n^2$ бол $p\mid n$ болохыг батал.}
\pr{$\sqrt2+\sqrt3$ иррационал болохыг батал. (Заавар: эсрэгээс нь $\sqrt2+\sqrt3=\frac ab$ гэж үз.)}
\pr{$\log_23$ иррационал болохыг батал.}
\pr{Анхны тоо хязгааргүй олон байдгийг батал.}
\pr{$a>b>0$, $n\in\N$ бол $a^n-b^n\ge n\,b^{n-1}(a-b)$ болохыг батал.}
\pr{Ньютоны бином $(a+b)^n=\sum_{k=0}^{n}\binom nk a^{n-k}b^k$, $\binom nk=\frac{n!}{k!(n-k)!}$ томьёог индукцээр батал.}
\pr{$4k+3$ хэлбэрийн анхны тоо хязгааргүй олон байдгийг батал. (Заавар: эсрэгээс нь $p_1,\dots,p_m$ гэж аваад $N=4p_1\cdots p_m-1$-ийг авч үз.)}
\end{plist}
\end{frame}
